\chapter{Estado de la cuestión} \label{chap-2-Estado-cuestion}
\section{Conceptos previos}
\textcolor{teal}{¿Introduzco un capítulo/sección 'Conceptos previos' que haga una introducción a los conceptos sobre los que se cimenta este trabajo? P.e.: Inteligencia Artificial, LLMs y otro tipo de conceptos de imprescindible conocimiento para la comprensión del trabajo.}

\textcolor{teal}{
    Qué incluir:
    \begin{itemize}
        \item Lenguaje natural y evaluación
        \item Qué es el lenguaje natural (texto y habla)
    \end{itemize}
    Qué implica “evaluar lenguaje”:
    \begin{itemize}
        \item forma (gramática, ortografía)
        \item contenido (semántica, adecuación)
        \item pragmática (intención, contexto)
        \item Modelos de lenguaje (LLMs)
        \item Definición de LLM
    \end{itemize} 
    Capacidades:
    \begin{itemize}
        \item generación
        \item comprensión
        \item evaluación (emergente)
        \item Idea clave:
            Los LLMs no fueron diseñados explícitamente para evaluar, pero muestran capacidad para hacerlo
    \end{itemize}
    ASR (Automatic Speech Recognition)
    \begin{itemize}
        \item Qué es
        \item Pipeline básico:
        \item audio → texto
    \end{itemize}
    Problemas:
    \begin{itemize}
        \item ruido
        \item acentos
        \item niños (importante para tu contexto)
    \end{itemize}
    Cierre del apartado:
    Introduce la idea de que evaluar lenguaje automático implica múltiples niveles → necesidad de enfoques complesjos
}

\section{Evaluación automática de lenguaje}
\textcolor{teal}{Revisión de técnicas de Automated Essay Scoring y Automatic Short Answer Grading. Revisión de los artículos académicos relacionados con el análisis y evaluación de lenguaje escrito.}

\textcolor{teal}{
    2.2.1 Automated Essay Scoring (AES)
    Qué explicar:
    \begin{itemize}
        \item Evaluación automática de textos largos
        \item Métodos tradicionales:
        \begin{itemize}
            \item features lingüísticas (longitud, vocabulario)
            \item modelos supervisados            
        \end{itemize}
    \end{itemize}
    Problemas:
    \begin{itemize}
        \item poca interpretabilidad
        \item dependencia de datasets anotados
        \item dificultad con contenido semántico
    \end{itemize}
    Transición a LLMs:
    \begin{itemize}
        \item LLMs permiten evaluar coherencia y generar feedback
    \end{itemize}
    Pero:
    \begin{itemize}
        \item riesgo de inconsistencia
        \item falta de control
    \end{itemize}
    2.2.2 Automatic Short Answer Grading (ASAG)
    Qué incluir:
    \begin{itemize}
        \item evaluación de respuestas cortas
        \item comparación con respuestas de referencia
    \end{itemize}
    Métodos:
    \begin{itemize}
        \item similitud semántica
        \item embeddings
        \item modelos supervisados
    \end{itemize}
    Problemas:
    \begin{itemize}
        \item múltiples respuestas correctas
        \item ambigüedad
    \end{itemize}
}
\section{Evaluación automática del habla}
\textcolor{teal}{Métodos de evaluación de pronunciación, fluidez y prosodia mediante técnicas de procesamiento del habla. Revisión de los artículos académicos y los sistemas existentes con el análisis y evaluación de las características del habla.}
\textcolor{teal}{
    Qué evaluar en habla:
    \begin{itemize}
        \item pronunciación
        \item fluidez
        \item ritmo (prosodia)
    \end{itemize}
    Métodos clásicos:
    \begin{itemize}
        \item análisis acústico
        \item modelos fonéticos
    \end{itemize}
    Problemas:
    \begin{itemize}
        \item alta variabilidad
        \item dependencia del idioma
        \item dificultad con niños (clave si lo quieres mencionar indirectamente)
    \end{itemize}
    Conexión con LLMs:
    \begin{itemize}
        \item LLM no trabaja directamente con audio
        \item depende de ASR
    \end{itemize}
    Idea clave -> la evaluación de habla se convierte en evaluación de texto + pérdida de información
}
\section{Sistemas híbridos y enfoques basados en recuperación de información}
\textcolor{teal}{Descripción de aproximaciones que combinan modelos lingüísticos, técnicas tradicionales y recuperación de información. Revisión de los artículos académicos orientados al uso de sistemas de evaluación co nrecuperación deinformación y recomendaciones.}

\textcolor{teal}{
    2.4.1 Sistemas híbridos
    Qué son:
    \begin{itemize}
        \item combinación de:
        \begin{itemize}
            \item reglas lingüísticas
            \item modelos clásicos
            \item LLMs
        \end{itemize}        
    \end{itemize}
    Ventajas:
    \begin{itemize}
        \item mayor control
        \item más interpretabilidad
        \item menor dependencia del modelo
    \end{itemize}    
    Ejemplo aplicado: corrector gramatical + LLM evaluador
}

\textcolor{teal}{
    2.4.2 RAG (Retrieval-Augmented Generation)
    Qué explicar:
    \begin{itemize}
        \item uso de información externa
        \item mejora de contexto
    \end{itemize}
    Aplicación en evaluación:
    \begin{itemize}
        \item enunciado
        \item solución modelo
        \item rúbrica
    \end{itemize}
    Ventajas:
    \begin{itemize}
        \item reduce alucinaciones
        \item mejora precisión
    \end{itemize}
    Problemas:
    \begin{itemize}
        \item complejidad
        \item dependencia de datos externos
    \end{itemize}
}

\textcolor{teal}{
    2.4.3 LLMs como evaluadores (LLM-as-a-judge)
    Qué explicar:
    \begin{itemize}
        \item uso de LLMs para evaluar texto
    \end{itemize}
    Problemas identificados en literatura:
    \begin{itemize}
        \item sesgo
        \item inconsistencia
        \item sensibilidad al prompt
    \end{itemize}
    Conexión con tu TFM -> Tu trabajo analiza precisamente estos problemas
}

\textcolor{teal}{
    CIERRE DEL ESTADO DEL ARTE
    Debes terminar con algo así 
    Hueco de investigación
    \begin{itemize}
        \item existen muchos trabajos en:
        \begin{itemize}
            \item AES
            \item ASAG
            \item LLMs
        \end{itemize}
    \end{itemize}    
    Pero, faltan enfoques que:
    \begin{itemize}
        \item combinen evaluación modular
        \item analicen consistencia
        \item integren texto + habla
        \item evalúen fiabilidad sin datasets grandes
    \end{itemize}
    Posicionamiento de tu TFM -> Este trabajo propone un marco modular basado en LLMs para la evaluación automática del lenguaje, analizando su fiabilidad, consistencia y comportamiento bajo distintas configuraciones.
}